% twn_fuse Linux driver 開發講義第二冊（v1.1，繁體中文）。建置：在本資料夾執行 latexmk，輸出 build/handout.pdf。
\documentclass[11pt,a4paper]{article}
\usepackage[top=20mm,bottom=20mm,left=20mm,right=20mm,headheight=14pt]{geometry}
\usepackage{fontspec}
\usepackage{xeCJK}
\setmainfont{Segoe UI}
\setsansfont{Segoe UI}
\setmonofont{Consolas}[Scale=0.92]
\setCJKmainfont{Microsoft JhengHei}[BoldFont={Microsoft JhengHei Bold}]
\setCJKsansfont{Microsoft JhengHei}[BoldFont={Microsoft JhengHei Bold}]
\setCJKmonofont{Microsoft JhengHei}
\xeCJKsetup{CJKecglue={\hskip 0.15em plus 0.05em}}
\linespread{1.3}
\setlength{\parindent}{0pt}
\setlength{\parskip}{0.55em}

\usepackage{xcolor}
\usepackage{graphicx}
\usepackage{booktabs}
\usepackage{tabularx}
\usepackage{xltabular}
\usepackage{array}
\usepackage{enumitem}
\setlist{itemsep=0.15em,topsep=0.25em,leftmargin=1.4em}
\usepackage[most]{tcolorbox}
\usepackage{listings}
\usepackage{tikz}
\usetikzlibrary{positioning,arrows.meta,fit,calc,backgrounds}
\usepackage{fancyhdr}
\usepackage[hidelinks]{hyperref}
\usepackage{titlesec}

\definecolor{ink}{HTML}{1F3A5F}
\definecolor{oldc}{HTML}{8A5A00}   \definecolor{oldbg}{HTML}{FFF6E5}
\definecolor{techc}{HTML}{1F5FA8}  \definecolor{techbg}{HTML}{EAF2FC}
\definecolor{gainc}{HTML}{2E7D32}  \definecolor{gainbg}{HTML}{EAF6EA}
\definecolor{qac}{HTML}{5B3F8C}    \definecolor{qabg}{HTML}{F3EEFA}
\definecolor{pitc}{HTML}{A23B2A}   \definecolor{pitbg}{HTML}{FCEDEA}
\definecolor{codebg}{HTML}{F7F7F5}
\definecolor{kern}{HTML}{DCE8F5}   \definecolor{user}{HTML}{F2EBDD}  \definecolor{hw}{HTML}{E3F1E3}

\titleformat{\section}{\Large\bfseries\color{ink}}{\thesection}{0.8em}{}
\titleformat{\subsection}{\large\bfseries\color{ink}}{\thesubsection}{0.6em}{}
\titlespacing*{\section}{0pt}{0.6em}{0.6em}
\titlespacing*{\subsection}{0pt}{0.9em}{0.35em}

\pagestyle{fancy}
\fancyhf{}
\fancyhead[L]{\small\color{gray}twn\_fuse driver v1.1 開發講義}
\fancyhead[R]{\small\color{gray}\leftmark}
\fancyfoot[C]{\small\thepage}
\renewcommand{\headrulewidth}{0.3pt}
\renewcommand{\sectionmark}[1]{\markboth{\thesection\ #1}{}}

% ---- 每一節固定的框 ----
\newtcolorbox{oneline}{colback=ink!6,colframe=ink,boxrule=0.6pt,arc=1mm,left=2mm,right=2mm,
  fontupper=\bfseries,before skip=0.4em,after skip=0.6em}
% 預設的「原本」是第一冊的 driver v1.0；原本是 v1.0.0 的 UIO 寫法時，寫 \begin{before}[標題]
\newtcolorbox{before}[1][原本怎麼做（driver v1.0）]{colback=oldbg,colframe=oldc,title={#1},
  fonttitle=\bfseries,boxrule=0.6pt,arc=1mm,breakable}
\newtcolorbox{tech}{colback=techbg,colframe=techc,title={用了什麼技術},fonttitle=\bfseries,
  boxrule=0.6pt,arc=1mm,breakable}
\newtcolorbox{gain}{colback=gainbg,colframe=gainc,title={改善了什麼},fonttitle=\bfseries,
  boxrule=0.6pt,arc=1mm,breakable}
\newtcolorbox{qa}[1]{colback=qabg,colframe=qac,title={面試會問（完整版，約 1 分鐘）：#1},
  fonttitle=\bfseries,boxrule=0.6pt,arc=1mm,breakable}
\newtcolorbox{qashort}[1]{colback=qabg,colframe=qac,title={#1},fonttitle=\bfseries,
  boxrule=0.6pt,arc=1mm,breakable}
\newtcolorbox{pit}{colback=pitbg,colframe=pitc,title={踩過的坑},fonttitle=\bfseries,
  boxrule=0.6pt,arc=1mm,breakable}

% ---- 程式碼：一律由 extract.py 從 repo 的實際檔案抽出，保留原始行號 ----
\lstset{language=C,basicstyle=\ttfamily\footnotesize,numbers=left,numberstyle=\tiny\color{gray},
  numbersep=6pt,xleftmargin=1.6em,
  backgroundcolor=\color{codebg},frame=single,rulecolor=\color{gray!40},breaklines=true,
  breakatwhitespace=true,columns=fullflexible,keepspaces=true,showstringspaces=false,tabsize=4,
  keywordstyle=\color{techc}\bfseries,commentstyle=\color{gray!90!black}\itshape,
  stringstyle=\color{oldc},captionpos=t,aboveskip=0.5em,belowskip=0.5em,
  morekeywords={u32,u64,size_t,bool,dma_addr_t,__u32,__u64,uint32_t,uint64_t}}
% \snip[語言]{標籤}：印出 build/snippets/<標籤>.c，標題是「檔案 : 起始行–結束行」
\newcommand{\snip}[2][C]{%
  \input{build/snippets/#2.tex}%
  \lstinputlisting[language=#1,firstnumber=\snipfirst,
    title={\scriptsize\ttfamily\snippath\ :\ \snipfirst--\sniplast}]{build/snippets/#2.c}}
% 板上與模擬的實際輸出：\boardfile[說明]{檔名}{起-迄行} 直接從 results/raw/linux_driver_v11_20261027/ 讀出原文，
% 不手抄；行首「# 」是腳本印出的指令本身
\lstdefinestyle{board}{language={},numbers=none,xleftmargin=0pt,
  basicstyle=\ttfamily\scriptsize\color{white},backgroundcolor=\color{black!82},
  rulecolor=\color{black!82}}
\newcommand{\logdir}{../../../results/raw/linux_driver_v11_20261027}
% 每個框都不超過 25 行，放進 minipage 讓它不會跨頁斷開
% 紀錄檔不存在時印出紅字「紀錄待補」，交付前不可以有
\newcommand{\boardfile}[3][]{%
  \IfFileExists{\logdir/#2}{%
  \par\noindent\begin{minipage}{\linewidth}%
  \lstinputlisting[style=board,linerange={#3},
    title={\scriptsize\ttfamily results/raw/linux\_driver\_v11\_20261027/\detokenize{#2}\ :\ #3\normalfont #1}]{\logdir/#2}%
  \end{minipage}\par}{\par\noindent\textcolor{red}{（紀錄待補：\detokenize{#2} #3）}\par}}
\lstnewenvironment{cmdline}{\lstset{language=bash,numbers=none,xleftmargin=0pt,
  basicstyle=\ttfamily\footnotesize}}{}

% \code{...}：行內程式碼，保留空白；在 _ / . 之後允許換行，長名稱才不會超出版面
\ExplSyntaxOn
\NewDocumentCommand \code { m }
  {
    \tl_set:Ne \l_tmpa_tl { \tl_to_str:n {#1} }
    \regex_replace_all:nnN { ([_/\.]) } { \1 \c{allowbreak} } \l_tmpa_tl
    \texttt { \tl_use:N \l_tmpa_tl }
  }
\ExplSyntaxOff
\emergencystretch=2em
\newcommand{\todo}[1]{\textcolor{red}{（待寫：#1）}}
\newcommand{\sref}[1]{\S\ref{#1}}
% \lineof{標籤}：該段程式碼在原始檔的起始行號（給速查表用，跟著程式碼自動更新）
\newcommand{\lineof}[1]{\begingroup\input{build/snippets/#1.tex}\snipfirst\endgroup}

\title{\bfseries twn\_fuse Linux driver 開發講義（第二冊）\\[0.3em]
  \large v1.1：中斷、開機載入、記憶體隔離與發佈}
\author{Wei-Jhe Huang}
\date{2026 年 10 月 27 日}
\renewcommand{\contentsname}{目錄}
\renewcommand{\figurename}{圖}
% 講義裡會重複出現的量測數字只寫在這裡一次，來源都是 results/raw/linux_driver_v11_20261027/。
% 改數字只改這個檔案。
%
% ---- 開機組 D：v1.0.0 的開機檔（BOOT.BIN 690d87cc5f4b6897、system.dtb 18fff66fcfb7798b），
%      v1.1 的 twn_fuse.ko（5f77eacd80e62272），driver 輪詢
% fps A/B，每組 6 次（兩輪、順序相反）：D1_ab_measure.log 第 28-31 行
\newcommand{\abUio}{228.94}
\newcommand{\abUioSix}{229.09}
\newcommand{\abDrvDev}{228.83}
\newcommand{\abDrvNc}{231.53}
% CPU 存取 DMA buffer 的頻寬（MB/s）：D2_bench.log 第 11 行（Normal NC）、第 20 行（Device）
\newcommand{\bwNcOne}{316}
\newcommand{\bwNcEight}{4,250}
\newcommand{\bwDevOne}{103}
\newcommand{\bwDevEight}{853}
\newcommand{\bwLoad}{81}
% ioctl 與 backbone 一幀的時間分解（輪詢、Normal NC）：D2_bench.log 第 7-10 行
\newcommand{\ioctlUs}{1.63}
\newcommand{\bbAppUs}{4168.3}
\newcommand{\bbDrvUs}{4164.8}
\newcommand{\bbPlUs}{4129.0}
\newcommand{\bbSysUs}{3.5}
\newcommand{\bbWaitUs}{35.8}
\newcommand{\bbCpuUs}{454.6}
\newcommand{\bbCpuPct}{10.9}
% pipeline 的分段時間：D6_stage_times.log 第 4-21 行
\newcommand{\splitNcIn}{0.09--0.10}
\newcommand{\splitDevIn}{0.14--0.15}
\newcommand{\splitNcPeriod}{4.31--4.33}
\newcommand{\splitDevPeriod}{4.35--4.38}


\begin{document}
\maketitle
\thispagestyle{empty}
\tableofcontents
\clearpage

\section*{導讀}
\addcontentsline{toc}{section}{導讀}
\markboth{導讀}{}
\todo{導讀}

\clearpage
\section{量測先行：driver 版本快的那 1\% 從哪裡來}
\label{sec:ab}

\begin{oneline}
第一冊量到 driver 版本比 UIO 版本快約 1\%，這次用四組實驗拆開，找到原因是 CPU 看 DMA buffer 的記憶體屬性從 Device 改成 Normal non-cacheable，讓每一幀「CPU 寫入輸入影像」那一步變快。
\end{oneline}

\begin{before}[原本怎麼做（第一冊第 10 節的結論）]
第一冊在同一顆 bitstream、同一次開機上比較，driver 版本平均快 0.8\%（另一次開機 1.0\%）。但 UIO 版本換成 driver 版本時，下面三件事是一起改的：
\begin{itemize}
\item \textbf{輪詢在哪裡做}：UIO 版本由 app 在 user space 讀 CTRL；driver 版本在 kernel 裡用 \code{readl_poll_timeout()}。
\item \textbf{buffer 在哪裡}：UIO 版本固定在 device tree 保留的 \code{0x6000_0000}；driver 版本從 CMA 配，這次在 \code{0x5010_0000}。
\item \textbf{CPU 端的記憶體屬性}：UIO 版本經 \code{/dev/mem} 映射（下面第 148--153 行），得到 Device-nGnRnE；driver 版本用 \code{dma_mmap_coherent()}，得到 Normal non-cacheable。
\end{itemize}
\snip{uio_devmem}
三件事一起變，所以第一冊只能寫「driver 版本沒有變慢」，不能說是哪一件讓它變快。
\end{before}

\begin{tech}
\begin{itemize}
\item \textbf{控制變因的 A/B 實驗}：四組設定（見下表）。第 1→2 組只換 buffer 在保留區裡的位置；第 2→3 組同時換了等硬體的地方（user space → kernel）與 buffer 的來源（保留區 → CMA）；第 3→4 組只換 CPU 端的記憶體屬性。每組跑 6 次，分兩輪、第二輪順序反過來，避免板子慢慢變熱或變慢時偏袒某一組。
\item \textbf{module parameter（模組參數）\code{map_device}}：載入 driver 時給 \code{map_device=1}，driver 改用 \code{pgprot_noncached()} 加 \code{remap_pfn_range()} 映射，得到和 \code{/dev/mem} 一樣的 Device 屬性。只給量測用。
\item \textbf{環境變數 \code{PL_ARENA_PA}}：讓 UIO 版本把 buffer 改放在保留區裡的另一個位置。
\item \textbf{\code{twn_bench} 第 3 部分}：直接量 CPU 寫入、讀取 buffer 的頻寬。
\item \textbf{\code{PIPE_TIMING=1}}：pipeline app 自己印出每一步花多少時間。
\end{itemize}
\end{tech}

\begin{gain}
\begin{itemize}
\item \textbf{找到原因}：只差記憶體屬性的兩組（第 3、4 組），fps 從 \abDrvDev{} 變成 \abDrvNc{}（+1.2\%）。前三組（\abUio{}、\abUioSix{}、\abDrvDev{}）相差不到 0.3 fps，所以輪詢位置與 buffer 位置合起來看不出影響；這兩件事在第 2→3 組是一起換的，沒有分開，也不能排除兩個效應正好互相抵銷。
\item \textbf{找到機制}：Normal non-cacheable 下，CPU 寫 buffer 的頻寬是 Device 的 3.1 倍（1-byte store）到 5.0 倍（8-byte store）；pipeline 每一幀「寫輸入」從 \splitDevIn{} ms 變成 \splitNcIn{} ms，差約 0.05 ms，和每幀週期的差（約 0.04 ms）相當（\splitDevPeriod{} 對 \splitNcPeriod{} ms）。
\item 第一冊留下的「分不出原因」可以結案：這 1\% 來自 driver 的映射方式，而且有機制上的解釋，不是量測誤差。
\end{itemize}
\end{gain}

\subsection{實際例子：四組實驗}

\begin{center}
\small
\begin{tabularx}{\textwidth}{l l l l X}
\toprule
組別 & 等硬體的地方 & buffer 位置 & CPU 端屬性 & 跟上一組只差 \\
\midrule
\code{uio} & app（user space） & \code{0x6000_0000} & Device-nGnRnE & （起點，v1.0.0 的做法） \\
\code{uio_68} & app（user space） & \code{0x6800_0000} & Device-nGnRnE & buffer 位置 \\
\code{drv_dev} & driver（kernel） & \code{0x5010_0000}（CMA） & Device-nGnRnE & 等硬體的地方與 buffer 來源（兩件一起換） \\
\code{drv_nc} & driver（kernel） & \code{0x5010_0000}（CMA） & Normal non-cacheable & 只有 CPU 端屬性 \\
\bottomrule
\end{tabularx}
\end{center}

每一組用 pipeline app 跑 120 幀（\code{FUSE_BB2HD=1}、head 在 PL），量 fps。實驗腳本的核心是 \code{run()}：先用第一冊的 \code{bind.sh} 把裝置交給對應的 driver，再跑三次：

\snip[bash]{ab_run}

\code{drv_dev} 這一組靠的是 driver 的 \code{map_device} 參數。\code{twn_mmap()} 在它設定時改走 \code{remap_pfn_range()}，屬性用 \code{pgprot_noncached()}，也就是 \code{/dev/mem} 給的 Device-nGnRnE；沒設定時走平常的 \code{dma_mmap_coherent()}：

\snip{drv_mmap}

\textbf{板上結果。}開機組 D（v1.0.0 的開機檔，v1.1 的 driver，這時 driver 是輪詢版）：

\boardfile[，第 27--31 行：每組 6 次的平均]{D1_ab_measure.log}{1-2,27-31}

\code{uio}、\code{uio_68}、\code{drv_dev} 三組都在 228.8 到 229.1 fps 之間，只有 \code{drv_nc} 跳到 \abDrvNc{} fps。各組 6 次的原始值在同一個檔案的第 3--26 行，組內的差距不超過 1 fps。

\textbf{為什麼記憶體屬性會影響 fps：直接量 CPU 存取 buffer 的速度。}\code{twn_bench} 第 3 部分在 buffer 的輸入區連續寫 1 MB、讀 1 MB：

\snip{bench_copy}

\boardfile[：第 11 行是預設的 Normal non-cacheable，第 20 行是 \code{map_device=1}]{D2_bench.log}{11-12,20}

寫入差很多（1-byte store \bwDevOne{} 對 \bwNcOne{} MB/s，8-byte store \bwDevEight{} 對 \bwNcEight{} MB/s），讀取一樣（都是 \bwLoad{} MB/s）。

\textbf{差異落在 pipeline 的哪一步。}pipeline app 開 \code{PIPE_TIMING=1}，會印出 producer（負責「寫輸入、啟動 PL、等 PL」的那個執行緒）最後一幀每一步的時間：

\boardfile[：前三次是 Normal non-cacheable，後三次是 \code{map_device=1}]{D6_stage_times.log}{3-21}

\code{write-in}（CPU 把這一幀的輸入影像寫進 buffer）從 \splitDevIn{} ms 變成 \splitNcIn{} ms，其他步驟（\code{backbone PL}、\code{head PL}）幾乎不變。每幀週期差約 0.04 ms，和 \code{write-in} 的差相當。寫輸入在 producer 的關鍵路徑上：寫完才能啟動 PL，所以它省下的時間直接變成 fps。

\subsection{原理}

\textbf{ARM 的兩種「不經過 cache」的記憶體。}arm64 的頁表每一頁都帶一個記憶體屬性。這次牽涉兩種，CPU 存取時都不經過 cache，差別在寫入能不能合併：

\begin{center}
\small
\begin{tabularx}{\textwidth}{l X X}
\toprule
 & Device-nGnRnE & Normal non-cacheable \\
\midrule
誰會得到 & \code{/dev/mem} 映射 no-map 區（第一冊第 3 節）；\code{pgprot_noncached()} & \code{dma_mmap_coherent()}（第一冊第 4 節）；\code{pgprot_dmacoherent()} \\
寫入 & 每一個 store 都是一筆獨立的匯流排交易，不合併（nG = non-Gathering）；寫入的回應必須由終點發出，中途的緩衝不能提早回應（nE = no Early write acknowledgement） & 連續的小 store 可以在 CPU 的 write buffer（CPU 送出寫入前暫放的小緩衝區）裡合併成一筆大的寫入 \\
讀取 & 每次都到 DDR & 每次都到 DDR（沒有 cache） \\
對齊 & 不能非對齊存取 & 可以 \\
\bottomrule
\end{tabularx}
\end{center}

所以寫入差 3 到 5 倍、讀取一樣，跟量到的完全對得上。讀取 81 MB/s 表示每次 8 bytes 的讀取大約要 100 ns，大約是一趟 DDR 讀取的延遲，兩種屬性都要付。

\textbf{為什麼只有 \code{write-in} 變快。}pipeline 的 CPU 對 buffer 做兩件事：寫輸入（store）、讀 heatmap（load）。store 受屬性影響、在關鍵路徑上；load 不受屬性影響，而且在另一個執行緒（consumer），和 PL 同時進行。所以只有寫輸入那一步的改善會反映到 fps。

\textbf{為什麼不乾脆用會 cache 的記憶體。}PL 經 HP port 存取 DDR，不會去看 CPU 的 cache（不是 coherent 的連接）。若 buffer 映射成 cacheable，CPU 寫完輸入後必須先把 cache 寫回 DDR，PL 寫完輸出後又必須先作廢 cache 再讀，driver 要提供 \code{dma_sync_single_for_device()} 這類同步介面，app 也要配合。這是可以做的下一步，這次沒有做。

\textbf{實驗設計的三個細節。}第一，兩輪的順序相反（\code{uio}→\code{drv_nc}，再 \code{drv_nc}→\code{uio}），板子若隨時間慢慢變快或變慢，兩輪的影響會互相抵銷。第二，\code{uio_68} 這一組是為了看 buffer 在 DDR 的位置會不會影響：同樣用 UIO，只換位置，fps 沒有變。第三，\code{uio_68}→\code{drv_dev} 一次換了兩件事（等硬體的地方、buffer 從保留區改成 CMA），要把它們也分開，還需要一組「driver 版、buffer 放在保留區」或「UIO 版、buffer 在 CMA」，這次沒有做；因為這三組的 fps 幾乎一樣，結論不受影響，但嚴格說只能講「兩者合起來看不出影響」。

\subsection{面試會問}

\begin{qa}{你說 driver 版本比較快，你怎麼知道是哪個改動造成的？}
driver 版本同時改了三件事：輪詢從 user space 搬到 kernel、buffer 從保留區改成 CMA、CPU 端的記憶體屬性從 Device 改成 Normal non-cacheable。我做了四組實驗，每組 6 次、兩輪順序相反：先只換 buffer 在保留區裡的位置，再換成 driver（輪詢位置和 buffer 來源一起換），最後只換記憶體屬性。前三組都在 229 fps 附近，只有最後一步從 \abDrvDev{} 跳到 \abDrvNc{}，所以原因是記憶體屬性；輪詢位置和 buffer 來源是一起換的，我只能說它們合起來看不出影響。再往下量機制：Device 屬性下每個 store 都是獨立的匯流排交易，Normal non-cacheable 可以在 write buffer 合併，CPU 寫 buffer 的頻寬差 3 到 5 倍；pipeline 裡「寫輸入」那一步從 0.14--0.15 ms 變成 0.09--0.10 ms，和每幀週期的差相當。讀取兩邊一樣慢，因為都不經過 cache。
\end{qa}

\subsection{關鍵函式速查}

\begin{center}
\small
\begin{tabularx}{\textwidth}{l l X}
\toprule
名稱 & 位置 & 做什麼 \\
\midrule
\code{map_device} & \code{twn_fuse.c}:\lineof{drv_param_map} & module parameter：映射成 Device 屬性（只給量測用） \\
\code{twn_mmap()} & \code{twn_fuse.c}:\lineof{drv_mmap} & 預設 \code{dma_mmap_coherent()}；\code{map_device} 時 \code{remap_pfn_range()} \\
\code{run()} & \code{test/ab_measure.sh}:\lineof{ab_run} & 切換 driver、跑 pipeline 三次 \\
第 3 部分 & \code{test/twn_bench.c}:\lineof{bench_copy} & CPU 寫、讀 buffer 的頻寬 \\
\code{v11_check.sh split} & \code{test/v11_check.sh}:\lineof{chk_split} & 兩種映射下的 pipeline 分段時間 \\
\bottomrule
\end{tabularx}
\end{center}

\clearpage
\section{ioctl 要花多少時間}
\label{sec:ioctl}

\begin{oneline}
一次 ioctl 的來回是 \ioctlUs{} µs。一幀 backbone 約 4.17 ms，其中系統呼叫和參數複製只占 \bbSysUs{} µs；看得見的額外成本是 driver 輪詢等硬體：多等 \bbWaitUs{} µs，而且 CPU 每幀花掉 \bbCpuUs{} µs。這個數字是第 3 節改用中斷的理由。
\end{oneline}

\begin{before}[原本怎麼做（第一冊沒有量）]
第一冊說 driver 版本「沒有變慢」，根據的是整條 pipeline 的 fps。但 ioctl 本身要多少時間、driver 等硬體的方式浪費多少 CPU，都沒有量。面試時被問「你的 driver 有多少 overhead」，只能回答「fps 沒掉」，說不出數字。
\end{before}

\begin{tech}
\begin{itemize}
\item \textbf{三個時鐘對照}：app 用 \code{clock_gettime(CLOCK_MONOTONIC)} 量 ioctl 前後；driver 在 ioctl 裡用 \code{ktime_get_ns()} 量 START 到 done，放進回傳的結構；PL 自己的 cycle counter 除以 240 MHz。三者的差就是各層的成本。
\item \textbf{\code{getrusage(RUSAGE_SELF)}}：這個 process 用掉的 CPU 時間（user 加 kernel），量輪詢佔用多少 CPU。
\item \textbf{空 ioctl 迴圈}：\code{TWN_IOC_STATUS} 不拿鎖、只讀六個暫存器，連打 100,000 次取平均。
\end{itemize}
\end{tech}

\begin{gain}
\begin{itemize}
\item ioctl 介面本身不是瓶頸：一次來回 \ioctlUs{} µs，一幀的系統呼叫加參數複製 \bbSysUs{} µs，不到幀時間的 0.1\%。
\item 找到真正的成本：輪詢的等待（driver 量的時間比 PL 的 cycle 多 \bbWaitUs{} µs）與 CPU 時間（每幀 \bbCpuUs{} µs，佔幀時間 \bbCpuPct{}\%）。第 3 節用中斷把它們降到 \irqWaitUs{} µs 與 \irqCpuUs{} µs。
\end{itemize}
\end{gain}

\subsection{實際例子}

\code{test/twn_bench.c} 第 1 部分：空的 ioctl 連打 100,000 次。

\snip{bench_status}

第 2 部分：同一幀反覆跑 1000 次，每次記下三個時間：

\snip{bench_bbrun}

\textbf{板上（開機組 D，driver 輪詢）：}

\boardfile{D2_bench.log}{6-10}

\begin{center}
\small
\begin{tabularx}{\textwidth}{l r X}
\toprule
量的是什麼 & 每幀 & 由哪兩個數字相減 \\
\midrule
app 看到的時間 & 4168.3 µs & ioctl 前後的 \code{clock_gettime} \\
driver 看到的時間 & 4164.8 µs & ioctl 裡 START 前到 done 後的 \code{ktime_get_ns} \\
PL 數的 cycle & 4129.0 µs & 991,000 cycles $\div$ 240 MHz \\
系統呼叫＋複製 & \bbSysUs{} µs & app $-$ driver \\
等待 & \bbWaitUs{} µs & driver $-$ PL \\
這個 process 的 CPU 時間 & \bbCpuUs{} µs & \code{getrusage} 的差除以 1000 \\
\bottomrule
\end{tabularx}
\end{center}

\subsection{原理}

\textbf{一次 ioctl 經過哪些地方。}app 呼叫 \code{ioctl()}，C 函式庫執行 \code{svc} 指令進入 kernel；kernel 的 \code{sys_ioctl} 用檔案描述子找到 \code{struct file}，呼叫 driver 的 \code{unlocked_ioctl}，也就是 \code{twn_ioctl()}；\code{twn_status()} 讀六個暫存器，\code{copy_to_user()} 把 32 bytes 的 \code{struct twn_status} 寫回，再從 kernel 返回。整條路 1.6 µs，其中六次暫存器讀取各要經過 AXI 到 PL 再回來。

\textbf{為什麼「等待」不全是 driver 的錯。}PL 的 cycle counter 只數 17 個 block 執行的時間，block 和 block 之間的空檔不算（第一冊第 8 節）。所以「driver $-$ PL」裡有一部分是 PL 自己的空檔；它是 driver 叫醒延遲的上限，不是精確值。第 3 節換成中斷後，這個差從 \bbWaitUs{} µs 降到 \irqWaitUs{} µs，降下來的部分才確定是輪詢造成的。

\textbf{為什麼輪詢會吃 CPU。}\code{readl_poll_timeout()} 每次睡 13 到 50 µs，用的是高精度計時器（hrtimer）。一幀 4.1 ms 要醒來上百次，每次都要經過計時器中斷、排程器、切回這個 process、讀一次暫存器、再睡。每次只花幾 µs，累積起來每幀約 0.45 ms。

\textbf{量測本身的誤差。}兩次 \code{clock_gettime} 本身也要時間，但它被平均到 100,000 次裡；\code{getrusage} 的解析度是 µs，總量是 0.45 秒等級，誤差可以忽略。同一個開機組裡，「等待」這一項在不同次執行之間會差 10 µs 左右（例如 \code{D2_bench.log} 第 9 行與第 18 行），所以只比較同一次開機、同一個程式的數字。

\subsection{面試會問}

\begin{qa}{你的 driver 有多少 overhead？}
我分三層量。空的 ioctl 一次來回 1.6 µs。跑一幀時，app 量到的時間只比 driver 在 kernel 裡量到的多 3.5 µs，這是系統呼叫加 \code{copy_from_user} 和 \code{copy_to_user}，不到幀時間的 0.1\%。driver 量到的又比 PL 自己數的 cycle 多 30 到 45 µs，這是輪詢的粒度加上 PL 在 block 之間的空檔；而且輪詢讓這個 process 每幀用掉 0.45 ms 的 CPU，大約 11\%。所以介面不是問題，問題是等待的方式。我改成中斷之後，CPU 時間降到每幀 \irqCpuUs{} µs，等待降到 \irqWaitUs{} µs。
\end{qa}

\subsection{關鍵函式速查}

\begin{center}
\small
\begin{tabularx}{\textwidth}{l l X}
\toprule
名稱 & 位置 & 做什麼 \\
\midrule
第 1 部分 & \code{test/twn_bench.c}:\lineof{bench_status} & 空 ioctl 連打 100,000 次 \\
第 2 部分 & \code{test/twn_bench.c}:\lineof{bench_bbrun} & app、driver、PL 三個時間與 CPU 時間 \\
\code{r.ns} & \code{twn_fuse.c} 的 \code{twn_bb_run()} & driver 在 kernel 裡量的 START 到 done \\
\bottomrule
\end{tabularx}
\end{center}

\clearpage
\section{中斷：從 RTL 的一根線到 driver 的 ISR}
\label{sec:irq}

\begin{oneline}
v1.0 的 bitstream 沒有中斷線，driver 只能每 13--50 µs 讀一次 done 位元。v1.1 在控制暫存器模組 \code{twn_csr} 加了一個 level 中斷輸出，一路接到 PS 的 \code{pl_ps_irq0}（GIC 的 SPI 89）；driver 改成睡在一個 completion 上，由中斷處理函式（ISR）叫醒。
\end{oneline}

\begin{before}
\textbf{硬體}：v1.0.0 的 block design 把 PS 的 PL 中斷輸入接成常數 0，只為了讓 PS 的設定不變：
\snip[tcl]{v10_irq_zero}
\textbf{driver v1.0}：等 done 只能輪詢（第一冊第 6 節）。\code{wait_bit()} 用 \code{readl_poll_timeout()}，每次讀完睡 13 到 50 µs：
\snip{v10_wait_bit}
\snip{v10_wait_done}
問題是一幀 4.1 ms 裡，CPU 要醒來讀暫存器上百次；而且 done 之後最多要再等一個睡眠週期才發現。
\end{before}

\begin{tech}
\begin{itemize}
\item \textbf{RTL}：新暫存器 IRQ（offset \code{0x0064}）放兩個致能位元；中斷輸出是暫存器輸出（registered）、level high。
\item \textbf{IP 封裝}：\code{ipx::infer_bus_interface} 把 \code{interrupt} 這個 port 宣告成 Vivado 的 interrupt 介面，參數 \code{SENSITIVITY = LEVEL_HIGH}。
\item \textbf{block design}：\code{connect_bd_net} 接到 \code{ps/pl_ps_irq0}。PetaLinux 讀 XSA 後自動在 device tree 節點寫 \code{interrupts = <0 89 4>}。
\item \textbf{kernel}：\code{platform_get_irq_optional()}、\code{devm_request_irq()}、ISR 回傳 \code{irqreturn_t}；\code{struct completion} 的 \code{init_completion()}、\code{reinit_completion()}、\code{complete()}、\code{wait_for_completion_timeout()}；module parameter \code{use_irq}。
\end{itemize}
\end{tech}

\begin{gain}
\begin{itemize}
\item \textbf{CPU 時間}：每幀從 \pollCpuUs{} µs 降到 \irqCpuUs{} µs（\code{twn_bench}，同一次開機、同一個程式，只差 \code{use_irq}）。
\item \textbf{等待延遲}：driver 量的時間減掉 PL 自己數的 cycle，從 \pollWaitUs{} µs 變成 \irqWaitUs{} µs。
\item \textbf{吞吐量不變差}：pipeline \fpsIrq{} fps（第 10 節的 board check）。
\item \textbf{相容}：device tree 沒有 \code{interrupts} 的舊 bitstream，driver 自動退回輪詢（開機組 D 的 dmesg 寫 \code{polling}）。
\end{itemize}
\end{gain}

\subsection{硬體：RTL 加一根線}

\code{twn_csr.v} 開頭的說明寫了整個設計：

\snip[Verilog]{csr_irq_doc}

\textbf{中斷怎麼產生。}每個 clock 都把「致能 AND done」重新算一次，存進暫存器 \code{irq}：

\snip[Verilog]{csr_irq_reg}

\code{ap_done} 是 backbone 做完、\code{hd_done_l} 是 head 做完的黏著位元（設起來後一直保持，直到被清掉）。所以 \code{irq} 在 done 設起來的下一個 cycle 拉高，在 done 被清掉的下一個 cycle 放下。中間的 \code{ifdef CSR_IRQ_BREAK_MASK} 是故意做壞的版本，只給模擬用，見下面的 gate。

\textbf{致能放在自己的暫存器。}寫 IRQ 暫存器設定兩個致能，讀它可以同時看到致能、兩個 done 和 \code{irq} 本身：

\snip[Verilog]{csr_irq_wr}
\snip[Verilog]{csr_irq_rd}

為什麼不沿用 CTRL 的 bit 2（\code{irq_en}）？因為 CTRL 的每一次寫入都會把 bit 2 一起寫掉：

\snip[Verilog]{csr_ctrl_wr}

driver 啟動一幀寫的是 \code{0x1}（START），清 done 寫的是 \code{0x2}；這兩個值的 bit 2 都是 0。若致能放在 CTRL[2]，每次啟動都會把中斷關掉。\code{CSR_IRQ_BREAK_CTRL} 就是這個錯誤的版本。

\textbf{單元 gate。}\code{rtl/fuse/sim/run_csr_irq.ps1} 跑三次：原設計必須 PASS，兩個故意做壞的版本必須 FAIL。後兩者證明 testbench 真的能抓到這兩種錯：

\boardfile[：PC 上的 Vivado xsim]{S1_csr_irq_gate.log}{2,4-7,9-15}

另外，整個 head 路徑的整合 gate（\code{tb_hd_int}）用改過的 \code{twn_csr.v} 重跑，4 個 operator、434,176 bytes 逐位元相同：

\boardfile{S2_hd_int_gate.log}{626,631,634}

\subsection{從 IP 到 device tree}

RTL 頂層 \code{twn_fuse_layer.v} 多一個 port，名字就叫 \code{interrupt}：

\snip[Verilog]{layer_irq_port}

封裝 IP 時把它宣告成 interrupt 介面。沒有這一步，Vivado 只把它當一條普通的線，PetaLinux 也不會在 device tree 寫出中斷：

\snip[tcl]{pkg_irq}

block design 把它接到 PS：

\snip[tcl]{bd_irq}

PetaLinux 從 XSA 產生的節點（開機組 E 的 \code{system.dtb}，用 \code{dtc} 反組譯）：

\boardfile[：PC 上的檢查，節點原文]{S3_binding_check.log}{24-33}

\code{interrupts = <0x00 0x59 0x04>} 的三個數字：0 表示 SPI（Shared Peripheral Interrupt，所有 CPU 共用的外部中斷），\code{0x59} = 89 是 SPI 編號，4 是 \code{IRQ_TYPE_LEVEL_HIGH}。GIC 內部的編號是 SPI 編號加 32，所以 \code{/proc/interrupts} 會顯示 121。

\subsection{driver：ISR 與 completion}

\textbf{probe 時申請中斷。}device tree 有中斷、且 \code{use_irq} 沒被關掉，就用 \code{devm_request_irq()} 註冊 ISR，再寫 IRQ 暫存器打開兩個致能；沒有中斷就留著輪詢：

\snip{drv_probe_irq}

\code{platform_get_irq_optional()} 和 \code{platform_get_irq()} 的差別是：找不到中斷時不在 kernel log 印錯誤，因為對舊 bitstream 來說這是正常的。

\textbf{ISR。}看哪一個 done 設起來，就清掉它、叫醒等它的人：

\snip{drv_isr}

\textbf{等待。}\code{wait_done()} 有中斷時睡在 completion 上，沒有時照舊輪詢：

\snip{drv_wait_done}

\textbf{啟動一幀。}和 v1.0 相比只多一行 \code{reinit_completion()}，而且位置有講究，見「原理」：

\snip{drv_bb_run}

head 的兩個函式也一樣：\code{twn_hd_start()} 在 START 之前 \code{reinit_completion(&td->hd_done)}，\code{twn_hd_wait()} 呼叫 \code{wait_done()}。

\textbf{卸載。}先關掉中斷來源，再讓 \code{devm} 釋放中斷：

\snip{drv_remove}

\textbf{板上：開機組 E。}driver 開機自動載入，說它用中斷；\code{/proc/interrupts} 有一條 GIC 121、Level：

\boardfile{E1_boot_state.log}{11-14}

\irqBoardEvidence

\subsection{原理}

\textbf{為什麼用 level，ISR 為什麼一定要清 done。}level 中斷的意思是「線拉高的期間，中斷一直成立」。GIC 在 ISR 結束（EOI）後會重新看這條線；如果 done 還沒清，線還是高，ISR 會立刻再被呼叫，形成中斷風暴（interrupt storm）。所以 ISR 一定要把中斷來源（done）清掉。level 的好處是不會漏：就算兩個 done 幾乎同時設起來，線會一直高到兩個都被處理。

\textbf{為什麼清完還要讀一次。}\code{writel()} 是 posted write：CPU 送出就繼續往下走，不等寫入真的到達 PL。ISR 若馬上返回，GIC 可能還看到舊的高電位，多進一次 ISR（這次兩個 done 都已清掉，會回傳 \code{IRQ_NONE}）。讀同一個裝置的暫存器時，Device 記憶體保證讀在寫之後到達，所以讀回來時，清除的寫入已經到了 PL。

\textbf{\code{reinit_completion()} 為什麼放在 START 之前。}completion 是一個計數加一個等待佇列：\code{complete()} 把計數加一並叫醒等待者，\code{wait_for_completion_timeout()} 在計數大於 0 時把它減一並返回。若 \code{reinit_completion()}（計數歸零）放在 START 之後，PL 若很快做完，ISR 的 \code{complete()} 可能先發生，接著被歸零，等待的人就睡到 3 秒逾時。放在 START 之前，PL 根本還沒開始，不可能有中斷搶先。

\textbf{為什麼 ISR 可以直接在硬體中斷裡做完。}它只做幾次暫存器讀寫和 \code{complete()}，都不會睡眠，可以在中斷的上下文（interrupt context）執行，不需要 threaded IRQ。

\textbf{為什麼致能要寫兩次（probe 打開、remove 關掉）。}中斷線在 FPGA 裡一直接著，driver 沒載入時致能是 reset 值 0，線不會動。probe 在 \code{devm_request_irq()} 之後才打開致能，remove 先關掉致能，再由 \code{devm} 釋放中斷；順序反過來，就可能有一個沒有人處理的中斷。

\subsection{面試會問}

\begin{qa}{你怎麼把輪詢改成中斷？從硬體講到軟體。}
硬體上，我在控制暫存器模組加了一個暫存器輸出的 level-high 中斷，等於「致能 AND done」；致能放在獨立的暫存器，因為 CTRL 每次寫入都會覆蓋原本的 bit 2，driver 一啟動就會把中斷關掉，這個錯誤我在 testbench 裡故意做出來、確認抓得到。IP 封裝時把 port 宣告成 interrupt 介面，block design 接到 PS 的 \code{pl_ps_irq0}，PetaLinux 就在 device tree 寫出 SPI 89、level high。driver 在 probe 用 \code{devm_request_irq()} 註冊 ISR；ISR 清掉 done、讀一次暫存器確保清除已經到達，再 \code{complete()}；等待的那一邊用 \code{wait_for_completion_timeout()}，\code{reinit_completion()} 放在 START 之前避免漏掉中斷。結果每幀 CPU 時間從 \pollCpuUs{} µs 降到 \irqCpuUs{} µs，fps 不變；device tree 沒有中斷時自動退回輪詢。
\end{qa}

\subsection{關鍵函式速查}

\begin{center}
\small
\begin{tabularx}{\textwidth}{l l X}
\toprule
名稱 & 位置 & 做什麼 \\
\midrule
\code{irq} 暫存器 & \code{rtl/fuse/twn_csr.v}:\lineof{csr_irq_reg} & 致能 AND done，level high \\
\code{A_IRQ} 寫入 & \code{rtl/fuse/twn_csr.v}:\lineof{csr_irq_wr} & 兩個致能 \\
interrupt 介面 & \code{vivado/package_fuse_ip.tcl}:\lineof{pkg_irq} & \code{SENSITIVITY = LEVEL_HIGH} \\
接到 PS & \code{vivado/build_bd_fuse.tcl}:\lineof{bd_irq} & \code{pl_ps_irq0} \\
\code{use_irq} & \code{twn_fuse.c}:\lineof{drv_param_irq} & 可以關掉中斷，退回輪詢 \\
\code{wait_done()} & \code{twn_fuse.c}:\lineof{drv_wait_done} & completion 或輪詢 \\
\code{twn_isr()} & \code{twn_fuse.c}:\lineof{drv_isr} & 清 done、讀回、\code{complete()} \\
probe 的中斷段 & \code{twn_fuse.c}:\lineof{drv_probe_irq} & \code{devm_request_irq()}、打開致能 \\
\bottomrule
\end{tabularx}
\end{center}

\clearpage
\section{重建 bitstream：timing closure}
\label{sec:timing}

\begin{oneline}
加中斷改了 RTL，bitstream 必須重建。用出貨時同一套建置配方，PL 時脈 240 MHz 下，第一次 route 就收斂：WNS +0.037 ns、WHS +0.008 ns。出貨的 bitstream 第一次 route 是 $-$0.030 ns，靠 post-route 的 \code{phys_opt_design} 才補到 0.000 ns。
\end{oneline}

\begin{before}[原本怎麼做（v1.0.0 的 bitstream）]
v1.0.0 的 bitstream（SHA-256 前 16 字元 \code{52b400785924efe0}）在 240 MHz 收斂得很勉強：第一次 route 後 WNS（Worst Negative Slack，最差路徑的時間餘裕；負的表示這條路徑來不及）是 $-$0.030 ns，建置腳本自動跑一次 post-route \code{phys_opt_design -directive AggressiveExplore}，補到 0.000 ns，完全沒有餘裕。出處：\code{results/board_results.json} 的 \code{timing_bitstream_52b400785924efe0}，報告原文在 \code{results/raw/thesis_model_20261012/timing_impl_fuse_ternpose.rpt}。

這個設計的 LUT 用了 93\%、block RAM 用了 97\%，改一點點 RTL，place \& route 的結果就可能不同。所以加中斷之前，最擔心的是 timing 收不了。
\end{before}

\begin{tech}
\begin{itemize}
\item \textbf{固定的建置配方}：\code{vivado/build.cmd} 裡的變數，包括 implementation strategy \code{Performance_NetDelay_high}、HP3 拆出獨立 interconnect、一段 block RAM 的位置限制（pblock）；這次一個都沒改。
\item \textbf{post-route 恢復}：\code{build_bd_fuse.tcl} 在 route 之後若 WNS < 0，依序跑三種 \code{phys_opt_design}，還不行就重新 route 再跑一次。
\item \textbf{資料路徑檢查}：合成後與 route 後各查一次，25 個資料路徑實例都還在 netlist 裡（沒被合成工具當成沒用的邏輯刪掉）。
\end{itemize}
\end{tech}

\begin{gain}
\begin{itemize}
\item timing 收斂，而且有正的餘裕：WNS +0.037 ns、WHS +0.008 ns，沒有用到 post-route 恢復。
\item 資源幾乎不變：LUT 多 24 個（44,081 對 44,057），block RAM 145 個、DSP 217 個都相同。
\item 這不代表加中斷讓 timing 變好；見「原理」。
\end{itemize}
\end{gain}

\subsection{實際例子}

\textbf{建置配方。}\code{build.cmd} 把配方寫成環境變數，\code{build_bd_fuse.tcl} 讀它們。這次只多設了一個 \code{FUSE_TAG=v11irq}，讓輸出檔名不蓋掉出貨版：

\snip[bash]{build_recipe}

\textbf{post-route 恢復。}以前出貨版就是靠這段補到 0.000 ns；這次 route 完 WNS 已經是正的，迴圈第一行就跳出：

\snip[tcl]{postroute}

\textbf{這次的建置紀錄}（\code{build.cmd} 的輸出，只留 timing 與檢查的行）：

\boardfile[：PC 上的 Vivado 2023.2]{S4_build_timing.log}{1,3-11}

各階段的 WNS：placement 後 $-$0.295 ns，placement 後的 \code{phys_opt} 把它拉到 $-$0.016 ns，route 後 +0.037 ns。第 8 行 \code{POSTROUTE start} 之後沒有 \code{phys_opt} 的紀錄，表示恢復步驟沒有執行。

\textbf{和出貨版比較}（兩邊都是 Vivado 2023.2、同一份配方、同一顆 xczu2cg-sfvc784-1-e，只差 RTL 多了中斷）：

\begin{center}
\small
\begin{tabularx}{\textwidth}{l r r X}
\toprule
 & v1.0.0（出貨） & v1.1（中斷） & 出處 \\
\midrule
第一次 route 的 WNS & $-$0.030 ns & +0.037 ns & \code{board_results.json}；\code{S4_build_timing.log} 第 7 行 \\
最後的 WNS & 0.000 ns & +0.037 ns & 兩份 \code{timing_impl} 報告第 142 行 \\
最後的 WHS & 0.002 ns & 0.008 ns & 同上 \\
post-route 恢復 & 用了一次 & 沒用到 & 報告第 10 行 Design State \\
CLB LUT & 44,057 & 44,081 & 兩份 \code{util_impl} 報告 \\
CLB Register & 36,339 & 36,253 & 同上 \\
Block RAM tile & 145 & 145 & 同上 \\
DSP & 217 & 217 & 同上 \\
\bottomrule
\end{tabularx}
\end{center}

\textbf{一個要記下的差異。}配方裡有一個 pblock，希望某一組 block RAM（\code{RB_T1}）放在第 0 到第 2 欄。出貨版有 1 顆落在範圍外，這次有 6 顆（紀錄第 10--11 行）。pblock 在這個配方裡只是建議性的位置限制，落在外面不算錯；timing 照樣收斂。若下一次重建 timing 收不了，這是第一個要看的地方。

\subsection{原理}

\textbf{為什麼加了邏輯，timing 反而比較好。}不是因為中斷。place \& route 是一個從某個起點出發的最佳化，netlist 改一點點（這次多 24 個 LUT），起點就不同，最後的結果像換了一個亂數種子。出貨版在同一份配方下也重建過：2026-10-11 從乾淨的 clone 重建，第一次 route 一樣是 $-$0.030 ns（\code{board_results.json} 的 \code{rebuilt_from_clone_20261011}）。所以 $-$0.030 和 +0.037 之間的差，是這種「換起點」的變化，不是設計變好。正確的說法是：這次收斂了，餘裕 +0.037 ns；同一份配方仍然沒有多少餘裕。

\textbf{中斷這條路徑本身很短。}\code{irq} 是一個暫存器，輸入是四個位元的 AND/OR，輸出直接到 PS。它不在關鍵路徑上；關鍵路徑仍然在運算陣列與 block RAM 之間（出貨版的報告與這次相同）。

\textbf{以前成功的作法。}這個專案過去收斂 timing 的經驗都寫在配方裡：固定的 strategy、把 HP3 拆出來減少 interconnect 的擁擠、限制 block RAM 的位置，以及 route 後自動的 \code{phys_opt}。這次的原則是「配方不動，只動 RTL」，這樣 timing 若出問題，原因只會是 RTL 的改動。

\subsection{面試會問}

\begin{qa}{改了 RTL 之後 timing 收不了，你會怎麼做？}
先確定建置配方沒變，這樣差異只來自 RTL。這個設計 LUT 用 93\%、block RAM 用 97\%，在 240 MHz 本來就沒什麼餘裕：出貨版第一次 route 是 $-$0.030 ns，靠 post-route 的 \code{phys_opt_design} 補到 0。我把這個恢復步驟寫進建置腳本：route 完若 WNS 小於 0，就依序跑三種 \code{phys_opt} directive，還不行再重新 route。這次加中斷，第一次 route 就是 +0.037 ns，恢復沒用到。我不會說中斷讓 timing 變好：同一份配方重建本來就有幾十 ps 的變化。還要確認合成沒有把資料路徑刪掉，所以建置腳本合成後和 route 後都檢查 25 個資料路徑實例還在。
\end{qa}

\subsection{關鍵函式速查}

\begin{center}
\small
\begin{tabularx}{\textwidth}{l l X}
\toprule
名稱 & 位置 & 做什麼 \\
\midrule
建置配方 & \code{vivado/build.cmd}:\lineof{build_recipe} & strategy、時脈、pblock 等變數 \\
post-route 恢復 & \code{vivado/build_bd_fuse.tcl}:\lineof{postroute} & WNS < 0 時跑 \code{phys_opt}，再不行重新 route \\
\code{FUSE_TAG} & \code{vivado/build.cmd} & 輸出檔名，不蓋掉出貨版 \\
\bottomrule
\end{tabularx}
\end{center}

\clearpage
\section{debugfs：除錯功能搬進 driver}
\label{sec:debugfs}
\todo{debugfs：除錯功能搬進 driver}

\clearpage
\section{開機自動載入：PetaLinux recipe}
\label{sec:autoload}
\todo{開機自動載入：PetaLinux recipe}

\clearpage
\section{拿掉 /dev/mem 與保留區}
\label{sec:devmem}
\todo{拿掉 /dev/mem 與保留區}

\clearpage
\section{SMMU：限制 PL 能碰的記憶體}
\label{sec:smmu}
\todo{SMMU：限制 PL 能碰的記憶體}

\clearpage
\section{device-tree binding 與靜態檢查}
\label{sec:binding}
\todo{device-tree binding 與靜態檢查}

\clearpage
\section{發佈 v1.1}
\label{sec:release}
\todo{發佈 v1.1}

\appendix
\clearpage
\section*{名詞表}
\addcontentsline{toc}{section}{名詞表}
\todo{名詞表}

\clearpage
\section{面試十題：30 秒版本}

每題先講結論，再講一個板上的數字或證據。完整版在括號中的章節。

\begin{qashort}{1. driver 版本比較快，是哪個改動造成的？（第 1 節）}
是 CPU 端的記憶體屬性。driver 用 \code{dma_mmap_coherent} 映射成 Normal non-cacheable，UIO 經 \code{/dev/mem} 是 Device-nGnRnE。四組 A/B 實驗裡，只有最後一步「只改屬性」有差：\abDrvDev{} 到 \abDrvNc{} fps；輪詢位置與 buffer 來源是一起換的，合起來看不出影響。機制是寫入可以合併，CPU 寫 buffer 快 3 到 5 倍，pipeline 每幀「寫輸入」從 0.14--0.15 ms 降到 0.09--0.10 ms。
\end{qashort}

\begin{qashort}{2. 你的 driver 有多少 overhead？（第 2 節）}
空 ioctl 一次 \ioctlUs{} µs；一幀的系統呼叫加參數複製 \bbSysUs{} µs，不到 0.1\%。真正的成本是輪詢：driver 比 PL 多等約 30 到 40 µs，CPU 每幀花 0.45 ms。同一次開機切換比較，改中斷後 CPU 從 \pollCpuUs{} µs 降到 \irqCpuUs{} µs、等待從 \pollWaitUs{} µs 降到 \irqWaitUs{} µs。
\end{qashort}

\begin{qashort}{3. 怎麼把輪詢改成中斷？（第 3 節）}
RTL 加一個暫存器輸出的 level 中斷，等於致能 AND done，致能放在獨立暫存器，免得被 START 寫掉；IP 宣告成 interrupt 介面，接 \code{pl_ps_irq0}，device tree 是 SPI 89。driver 用 \code{devm_request_irq}，ISR 清 done、讀回確保清除到達、\code{complete()}；\code{reinit_completion} 放在 START 之前。1000 幀剛好 1000 次中斷；同一次開機交替量，pipeline 從 \hPollFps{} 變成 \hIrqFps{} fps。
\end{qashort}

\begin{qashort}{4. 改了 RTL 之後 timing 收不了怎麼辦？（第 4 節）}
先固定配方，只動 RTL。建置腳本在 route 後若 WNS 小於 0，自動跑 \code{phys_opt_design}，不行再重新 route。這次把 AxPROT[1] 從 0 改成 1，只改常數，route 完卻是 $-$0.032 ns，靠這一步補到 0。我也準備了備案：Vivado 對同樣的輸入、到 placement 為止給出同樣的結果，打算重現上一版收斂的設計，再用 ECO 只改那幾根常數 pin；最後沒有用到。
\end{qashort}

\begin{qashort}{5. 除錯功能放哪裡？（第 5 節）}
放 debugfs：\code{debugfs_create_regset32} 一次印出暫存器，\code{DEFINE_DEBUGFS_ATTRIBUTE} 做可寫的 \code{nblk}，寫入時檢查範圍、拿 backbone 的鎖。它不是 ABI，app 不該依賴；寫 99 會回 \code{EINVAL}。
\end{qashort}

\begin{qashort}{6. driver 怎麼開機自動載入？（第 6 節）}
Yocto 的 module recipe：\code{inherit module} 用映像自己的 kernel 樹編，vermagic 一定對；\code{KERNEL_MODULE_AUTOLOAD} 讓 systemd 開機載入；device tree 把 \code{compatible} 改成 driver 認得的字串，不再需要 \code{driver_override}。板上開機第 4.8 秒就綁好。
\end{qashort}

\begin{qashort}{7. 為什麼限制 /dev/mem？（第 7 節）}
它把整個實體位址開給 root。driver 用 DMA API 配 buffer、只 mmap buffer，所以 app 不再需要 \code{/dev/mem}，256 MB 保留區也拿掉了。\code{STRICT_DEVMEM} 擋 kernel 管理的 System RAM，\code{IO_STRICT_DEVMEM} 擋 driver 佔用中的暫存器；板上 \code{devmem} 都回 \code{Operation not permitted}。邊界也量過：卸載 driver 後暫存器又讀得到，System RAM 以外的 DDR 最上面 1 MB 也讀得到。
\end{qashort}

\begin{qashort}{8. SMMU 怎麼讓 PL 只能碰自己的 buffer？（第 8 節）}
device tree 打開 SMMU，給加速器的四個 HP 埠登記 stream ID；DMA API 就自動把 buffer 映射到 IOVA，帶著這些 ID 的 Non-secure 交易只能存取映射過的頁，越界讀取會變成 context fault。第一次失敗是因為 AXI master 送的是 Secure 交易，Linux 只設定 Non-secure 那一側；改成 AxPROT[1] = 1 才生效。限制：發佈的 kernel 對沒登記的 stream ID 是 bypass，非 0 的 AXI ID 或 Secure 交易都擋不住。
\end{qashort}

\begin{qashort}{9. 怎麼確定 device tree 節點寫對了？（第 9 節）}
寫 dt-schema 的 binding，規定 v1.1 一定要有中斷、v1.0 不准有。官方檢查要裝 dtschema，這次環境不能裝，所以用自寫的部分檢查比對 driver 的比對表、編譯範例、檢查 PetaLinux 產生的實際節點，因此發現要描述 \code{interrupt-names}。
\end{qashort}

\begin{qashort}{10. 發佈映像怎麼處理密碼？（第 10 節）}
不放真正的密碼：用公開的預設密碼，建置時就用 Yocto 的 \code{passwd-expire} 讓它過期，所以序列埠與 ssh 的第一次登入都要先改；改之前 ssh 帶指令和 scp 都被拒絕，序列埠上 Ctrl-C、空密碼、兩次不一致都拿不到 shell。第一版用 \code{profile.d} 腳本，擋不住非互動的 ssh，才改掉。fixture 換成公式產生的合成圖，所以發佈檔沒有個人資料。限制是 initramfs，重開機就回到過期的預設密碼。
\end{qashort}

\clearpage
\section*{如何重現}
\addcontentsline{toc}{section}{如何重現}
\todo{如何重現}

\end{document}
